% Partie C du rapport : optimisation et convergence.
% Pour fusionner avec le rapport commun, ne garder que ce qui se trouve
% entre \begin{document} et \end{document} (ou utiliser \input).
% Les figures sont lues dans resultats/experiences/.

\documentclass[12pt,a4paper]{article}
\usepackage[utf8]{inputenc}
\usepackage[T1]{fontenc}
\usepackage[french]{babel}
\usepackage{amsmath,amssymb}
\usepackage{graphicx}
\usepackage{booktabs}
\usepackage[margin=2.5cm]{geometry}

\graphicspath{{../resultats/experiences/}}

\begin{document}

\section{Optimisation et convergence}

\subsection{Objectif de l'optimisation et fonction de coût}

Notre réseau a la forme $57 \to 32 \to 10$. Ses paramètres sont les
matrices de poids $W_1$ et $W_2$ et les biais $b_1$ et $b_2$, soit en tout
$57 \times 32 + 32 + 32 \times 10 + 10 = 2186$ paramètres, que l'on
regroupe dans un vecteur $\theta$.

Pour un lot de $N$ morceaux et $K = 10$ genres, le coût est l'erreur
quadratique moyenne
\[
E(\theta) = \frac{1}{NK} \sum_{n=1}^{N} \sum_{k=1}^{K}
\bigl(a_{n,k} - y_{n,k}\bigr)^2 ,
\]
où $a_{n,k}$ est la sortie sigmoïde du réseau et $y_{n,k}$ la valeur
attendue (encodage one-hot). Apprendre revient donc à résoudre
$\min_{\theta} E(\theta)$. Le gradient $\nabla E(\theta)$ est calculé par
la rétropropagation, qui est la partie du membre B. Notre travail ici
consiste à utiliser ce gradient pour modifier les paramètres.

\subsection{Descente de gradient}

Pour un déplacement $d$ petit, un développement à l'ordre un donne
\[
E(\theta + d) \approx E(\theta) + \nabla E(\theta)^{T} d .
\]
Si on choisit $d = -\eta \, \nabla E(\theta)$ avec $\eta > 0$, on obtient
\[
E(\theta + d) \approx E(\theta) - \eta \, \|\nabla E(\theta)\|^2 ,
\]
donc le coût diminue tant que $\eta$ est assez petit et que le gradient
n'est pas nul. D'où la règle de mise à jour
\[
\theta_{t+1} = \theta_t - \eta \, \nabla E(\theta_t) .
\]
Cette approximation n'est valable que pour un pas petit. Si $\eta$ est trop
grand, le terme négligé peut dominer et le coût peut augmenter, ce qui
explique l'importance du taux d'apprentissage.

Dans notre code, le gradient est calculé sur des mini-lots de 64 morceaux
tirés au hasard à chaque époque. Le jeu d'apprentissage compte 787
morceaux, ce qui fait 13 mises à jour par époque. Le gradient d'un
mini-lot est une estimation bruitée du gradient complet. Une époque est
un passage sur tout le jeu d'apprentissage.

\subsection{Choix des taux d'apprentissage et protocole expérimental}

Pour que les comparaisons soient justes, toutes les expériences respectent
les mêmes conditions :
\begin{itemize}
  \item mêmes données (celles préparées par le membre A) ;
  \item 800 époques et mini-lots de 64 morceaux ;
  \item même graine aléatoire (42), donc mêmes poids de départ et même
        ordre de tirage des mini-lots ;
  \item un seul paramètre change d'une expérience à l'autre ;
  \item le coût est mesuré à la fin de chaque époque sur tout le jeu
        d'apprentissage.
\end{itemize}

Avant les expériences, nous avons vérifié que notre boucle reproduit celle
du membre B : avec $\eta = 0{,}5$, sans momentum, nous retrouvons son coût
final (0,0308) et ses précisions (84,75~\% en apprentissage et 72,22~\% en
test).

Les expériences réalisées sont les suivantes.

\begin{center}
\begin{tabular}{lll}
\toprule
Expérience & Paramètre qui varie & Valeurs \\
\midrule
1. Taux d'apprentissage & $\eta$ & 0,001 ; 0,01 ; 0,1 ; 0,5 ; 1,0 \\
2. Momentum & $\beta$ (avec $\eta = 0{,}1$) & 0 ; 0,5 ; 0,9 \\
3. Taille du réseau & couche cachée ($\eta = 0{,}5$) & 8 ; 16 ; 32 ; 64 ; 128 \\
\bottomrule
\end{tabular}
\end{center}

Le temps de calcul est mesuré pour chaque essai. La précision sur le jeu
de test est donnée à titre indicatif : l'évaluation officielle
(précision et matrice de confusion) est la partie du membre A. Le jeu de
test contient 198 morceaux, donc un morceau correspond à environ
0,5 point de précision.

\subsection{Le momentum}

Le momentum garde en mémoire une partie de la mise à jour précédente.
Avec une vitesse $v$ initialisée à zéro, on utilise
\[
v_t = \beta \, v_{t-1} + \nabla E(\theta_t), \qquad
\theta_{t+1} = \theta_t - \eta \, v_t .
\]
Avec $\beta = 0$, on retrouve exactement la descente de gradient simple.

Si le gradient reste à peu près constant, la vitesse tend vers
$v = g/(1-\beta)$, donc le pas effectif devient $\eta/(1-\beta)$. Pour
$\beta = 0{,}9$, le pas est environ dix fois plus grand dans une direction
stable. Il faut en tenir compte quand on compare avec la descente simple.

Le code du membre B écrit le momentum sous la forme
$v' \leftarrow \mu v' - \eta g$ puis $w \leftarrow w + v'$. C'est la même
méthode : il suffit de poser $v' = -\eta v$ et $\mu = \beta$.

\subsection{Résultats et analyse de convergence}

\subsubsection{Influence du taux d'apprentissage}

La figure~\ref{fig:exp1} montre le coût en fonction de l'époque pour les
cinq valeurs de $\eta$, et le tableau~\ref{tab:exp1} donne les valeurs
après 800 époques.

\begin{figure}[ht]
\centering
\includegraphics[width=0.8\textwidth]{exp1_taux_apprentissage.png}
\caption{Coût en fonction de l'époque pour plusieurs valeurs de $\eta$
(descente simple).}
\label{fig:exp1}
\end{figure}

\begin{table}[ht]
\centering
\begin{tabular}{ccccc}
\toprule
$\eta$ & Coût final & Préc. train (\%) & Préc. test (\%) & Temps (s) \\
\midrule
0,001 & 0,1118 & 13,98 & 14,65 & 4,0 \\
0,01  & 0,0834 & 33,80 & 33,84 & 3,7 \\
0,1   & 0,0599 & 58,58 & 57,58 & 3,7 \\
0,5   & 0,0308 & 84,75 & 72,22 & 4,0 \\
1,0   & 0,0199 & 91,87 & 72,73 & 3,8 \\
\bottomrule
\end{tabular}
\caption{Résultats de l'expérience 1 après 800 époques.}
\label{tab:exp1}
\end{table}

Plus $\eta$ est grand, plus le coût baisse vite : le coût final passe de
0,1118 pour $\eta = 0{,}001$ à 0,0199 pour $\eta = 1{,}0$. Avec
$\eta = 0{,}001$, le réseau apprend très lentement : après 800 époques sa
précision d'apprentissage est de 14~\%, alors qu'un choix au hasard donnerait
environ 10~\% avec dix genres. Les courbes sont régulières. Nous n'avons
observé ni oscillation ni divergence, même pour $\eta = 1{,}0$ : dans
l'intervalle testé, nous n'avons donc pas atteint de taux d'apprentissage
trop grand.

Un coût plus bas n'apporte pas autant en test. En passant de
$\eta = 0{,}5$ à $\eta = 1{,}0$, le coût baisse de 0,0308 à 0,0199 et la
précision d'apprentissage monte de 84,75~\% à 91,87~\%, mais la précision de
test passe seulement de 72,22~\% à 72,73~\%, soit un morceau de plus sur 198.
L'écart entre apprentissage et test augmente (environ 12,5 points puis
19,1 points), ce qui indique un début de surapprentissage.

\subsubsection{Descente simple et momentum}

Les résultats sont donnés par la figure~\ref{fig:exp2} et le
tableau~\ref{tab:exp2}. Le pas $\eta$ vaut 0,1 pour les trois courbes.

\begin{figure}[ht]
\centering
\includegraphics[width=0.8\textwidth]{exp2_momentum.png}
\caption{Descente simple et momentum, avec $\eta = 0{,}1$.}
\label{fig:exp2}
\end{figure}

\begin{table}[ht]
\centering
\begin{tabular}{cccccc}
\toprule
$\beta$ & Pas effectif & Coût final & Préc. train (\%) & Préc. test (\%)
 & Temps (s) \\
\midrule
0   & 0,1 & 0,0599 & 58,58 & 57,58 & 4,0 \\
0,5 & 0,2 & 0,0481 & 72,68 & 62,63 & 4,0 \\
0,9 & 1,0 & 0,0200 & 92,25 & 73,23 & 4,3 \\
\bottomrule
\end{tabular}
\caption{Résultats de l'expérience 2 après 800 époques ($\eta = 0{,}1$).}
\label{tab:exp2}
\end{table}

Le momentum fait baisser le coût beaucoup plus vite : à $\eta = 0{,}1$, le
coût final passe de 0,0599 sans momentum à 0,0481 pour $\beta = 0{,}5$ et à
0,0200 pour $\beta = 0{,}9$. Cela correspond à la remarque sur le pas
effectif $\eta/(1-\beta)$, qui vaut 0,2 pour $\beta = 0{,}5$ et 1,0 pour
$\beta = 0{,}9$. La courbe $\beta = 0{,}9$ est presque la même que celle de
la descente simple avec $\eta = 1{,}0$ (coûts finaux 0,01997 et 0,01995,
précisions d'apprentissage 92,25~\% et 91,87~\%). Dans nos essais, le
momentum se comporte donc surtout comme un pas plus grand. Nous n'avons pas
mis en évidence d'autre avantage, par exemple une meilleure stabilité,
puisque la descente simple n'oscille pas non plus pour $\eta = 1{,}0$. La
précision de test pour $\beta = 0{,}9$ (73,23~\%) ne diffère que d'un morceau de
celle obtenue avec $\eta = 1{,}0$ (72,73~\%).

\subsubsection{Influence de la taille du réseau}

Les résultats sont donnés par la figure~\ref{fig:exp3} et le
tableau~\ref{tab:exp3}, avec $\eta = 0{,}5$ et sans momentum. Le nombre de
paramètres vaut $68h + 10$ pour $h$ neurones cachés.

\begin{figure}[ht]
\centering
\includegraphics[width=0.8\textwidth]{exp3_taille_reseau.png}
\caption{Coût en fonction de l'époque pour plusieurs tailles de couche
cachée.}
\label{fig:exp3}
\end{figure}

\begin{table}[ht]
\centering
\begin{tabular}{cccccc}
\toprule
Neurones cachés & Paramètres & Coût final & Préc. train (\%)
 & Préc. test (\%) & Temps (s) \\
\midrule
8   & 554  & 0,0438 & 74,08 & 63,13 & 2,0 \\
16  & 1098 & 0,0353 & 83,35 & 70,20 & 2,9 \\
32  & 2186 & 0,0308 & 84,75 & 72,22 & 3,9 \\
64  & 4362 & 0,0288 & 86,79 & 71,21 & 8,8 \\
128 & 8714 & 0,0282 & 86,66 & 72,73 & 16,1 \\
\bottomrule
\end{tabular}
\caption{Résultats de l'expérience 3 après 800 époques. Les temps dépendent
de l'ordinateur utilisé.}
\label{tab:exp3}
\end{table}

Entre 8 et 32 neurones, le coût et la précision s'améliorent nettement
(précision de test de 63,13~\% à 72,22~\%). Au-delà de 32 neurones, le gain sur
le coût est faible (0,0308, 0,0288 puis 0,0282), alors que le temps de
calcul augmente fortement (de 3,9~s à 16,1~s) et que le nombre de paramètres
double à chaque fois. Les précisions de test pour 32, 64 et 128 neurones
(72,22~\%, 71,21~\% et 72,73~\%) ne diffèrent que de un à trois morceaux : on ne
peut donc pas conclure qu'un réseau plus grand est meilleur en test.
Une taille de 32 neurones nous semble un bon compromis entre coût et temps de
calcul.

\subsubsection{Convergence}

Après 800 époques, aucune des courbes n'est stabilisée : elles baissent
encore, y compris celle de $\eta = 1{,}0$. La baisse est très lente pour
$\eta = 0{,}01$, mais elle continue. On ne peut donc pas dire que le réseau
a convergé, seulement que le coût diminue régulièrement. Les courbes ne
commencent pas toutes à la même hauteur parce que le coût est mesuré à la
fin de la première époque, après 13 mises à jour : un grand $\eta$ a déjà fait
baisser le coût à ce moment-là.

\subsection{Tableau de comparaison des méthodes}

\begin{table}[ht]
\centering
\small
\setlength{\tabcolsep}{4pt}
\begin{tabular}{lccccccc}
\toprule
Méthode & $\eta$ & $\beta$ & Pas eff. & Coût final & Train (\%) & Test (\%)
 & Temps (s) \\
\midrule
Descente simple & 0,1 & 0   & 0,1 & 0,0599 & 58,58 & 57,58 & 3,7 \\
Descente simple & 0,5 & 0   & 0,5 & 0,0308 & 84,75 & 72,22 & 4,0 \\
Descente simple & 1,0 & 0   & 1,0 & 0,0199 & 91,87 & 72,73 & 3,8 \\
Momentum        & 0,1 & 0,5 & 0,2 & 0,0481 & 72,68 & 62,63 & 4,0 \\
Momentum        & 0,1 & 0,9 & 1,0 & 0,0200 & 92,25 & 73,23 & 4,3 \\
\bottomrule
\end{tabular}
\caption{Comparaison des méthodes après 800 époques (réseau 57-32-10).}
\label{tab:comparaison}
\end{table}

\subsection{Limites et conclusion}

Nos conclusions ont plusieurs limites. Chaque essai n'est fait qu'avec une
seule graine, donc une différence faible entre deux méthodes peut venir
du hasard. Le jeu de test est petit (198 morceaux), ce qui rend les écarts
de test de quelques morceaux peu significatifs. Les 800 époques ne suffisent
pas pour voir une stabilisation du coût. Nous n'avons pas cherché la valeur
de $\eta$ à partir de laquelle l'apprentissage devient instable. Les
mini-lots ajoutent du bruit au gradient. Enfin, le réseau utilise des
sigmoïdes avec une erreur quadratique, ce qui n'est pas le choix le plus
classique pour de la classification (softmax et entropie croisée sont une
extension possible).

Nous n'avons pas testé les méthodes d'ordre deux. La méthode de Newton
utilise la hessienne : $\theta_{t+1} = \theta_t - H^{-1}\nabla E(\theta_t)$.
Ici $H$ serait une matrice $2186 \times 2186$, soit environ 4,8 millions
d'entrées, et résoudre le système coûte de l'ordre de $2186^3
\approx 10^{10}$ opérations par itération. De plus, la hessienne n'est
pas forcément définie positive pour un réseau de neurones. C'est une
raison de rester sur des méthodes d'ordre un.

En résumé, dans nos expériences, le coût baisse d'autant plus vite que le pas
effectif est grand, que ce soit en augmentant $\eta$ ou en ajoutant du
momentum. Le momentum avec $\eta = 0{,}1$ et $\beta = 0{,}9$ donne presque le
même résultat que la descente simple avec $\eta = 1{,}0$. Un coût plus bas ne
donne pas une meilleure précision de test : au-delà de $\eta = 0{,}5$ et de
32 neurones cachés, la précision de test ne change pas de façon nette, alors
que l'écart avec la précision d'apprentissage grandit. Le réglage le plus
intéressant parmi ceux testés est donc un réseau de 32 neurones avec un pas
effectif de l'ordre de 0,5 à 1.

\end{document}
